% =================================================
% Материал для учителя
% =================================================

\packtitle{Методический комментарий}{Скрытый треугольник и центр описанной окружности}

\begin{teacherbox}[Назначение комплекта]
Комплект построен вокруг одной составной задачи. Его цель — не только получить
ответ, но и научить ученика выполнять цепочку распознаваний:
\[
\begin{aligned}
\text{достроение}
&\longrightarrow
\text{средняя линия}
\longrightarrow
\text{серединные перпендикуляры}\\
&\longrightarrow
\text{центр окружности}
\longrightarrow
\text{работа с углами}.
\end{aligned}
\]
Подготовительные задания соответствуют конкретным переходам основного решения,
поэтому их не следует рассматривать как независимую подборку упражнений.
\end{teacherbox}

\section{Дидактическая карта}

\begin{tabularx}{\linewidth}{@{}>{\bfseries\sffamily\raggedright\arraybackslash}p{42mm}>{\raggedright\arraybackslash}X@{}}
\toprule
Целевой результат & Ученик самостоятельно замечает полезное достроение,
обосновывает появление центра описанной окружности и находит искомый угол.\\
\midrule
Предварительные знания & Средняя линия треугольника; подобие треугольников;
серединный перпендикуляр; центр описанной окружности; центральный и вписанный
углы; прямоугольный равнобедренный треугольник.\\
\midrule
Главное затруднение & Переход от исходной трапеции к треугольнику \(BNC\).
Без этого построения прямые углы при \(A\) и \(D\) не раскрывают своей роли.\\
\midrule
Культура доказательства & Ученик различает «перпендикуляр» и
«серединный перпендикуляр», не делает вывод о коллинеарности по рисунку и явно
указывает общую дугу для центрального и вписанного углов.\\
\bottomrule
\end{tabularx}

\begin{routebox}[Карта полного решения]
\[
AD=\frac12BC,
\ AD\parallel BC
\Longrightarrow
A,D\text{ — середины }BN,CN,
\]
\[
AM\perp BN,
\ DM\perp CN
\Longrightarrow
M\text{ — центр описанной окружности }\triangle BNC,
\]
\[
ME\perp BC,
\ ME=AD=\frac12BC
\Longrightarrow
\angle MBC=45^\circ,
\]
\[
\angle BCN=64^\circ
\Longrightarrow
\angle BMN=128^\circ
\Longrightarrow
\angle MBN=26^\circ,
\]
\[
\angle ABC=45^\circ+26^\circ=71^\circ.
\]
\end{routebox}

\section{Рекомендуемый сценарий работы}

\subsection{Вариант 1. Два занятия по 45 минут}

\begin{tabularx}{\linewidth}{@{}p{19mm}p{45mm}X@{}}
\toprule
Время & Этап & Действия учителя\\
\midrule
5 мин & Постановка задачи & Показать только исходный рисунок. Не называть
точку \(N\) и не упоминать окружность.\\
10 мин & Поиск достроения & Вести вопросами: «Где встречается половина
параллельной стороны?», «Можно ли получить треугольник?».\\
12 мин & Средняя линия & Разобрать задания A–B; добиться явного вывода,
что \(A,D\) — середины сторон.\\
12 мин & Серединные перпендикуляры & Разобрать задания C–D и завершить пункт а).\\
6 мин & Рефлексия & Ученик формулирует, какую скрытую фигуру удалось найти.\\
\midrule
8 мин & Актуализация & Восстановить цепочку пункта а) без текста.\\
12 мин & Угол \(45^\circ\) & Ввести точку \(E\), разобрать хорду и треугольник
\(BME\).\\
12 мин & Угол \(26^\circ\) & Связать \(\angle BCN\) с \(\angle BMN\), затем
использовать равнобедренный треугольник.\\
7 мин & Сборка & Получить \(71^\circ\), записать полное решение.\\
6 мин & Ошибка и обобщение & Обсудить ложную коллинеарность \(N,M,E\) и
формулу \(135^\circ-\alpha\).\\
\bottomrule
\end{tabularx}

\subsection{Вариант 2. Одно занятие сильной группы}

\begin{teacherbox}[Что можно сократить]
Задания A–D можно дать как устную цепочку по одному рисунку. Задания E–G
оставить письменно, поскольку именно в них ученики чаще путают основания для
выводов. Полное решение после обсуждения записывается самостоятельно.
\end{teacherbox}

\section{Цепочка наводящих вопросов}

\begin{enumerate}
  \item Какую фигуру можно получить, если продолжить \(AB\) и \(CD\)?
  \item Как связаны треугольники \(NAD\) и \(NBC\)?
  \item Что означает коэффициент подобия \(\frac12\) для точек \(A\) и \(D\)?
  \item К какой прямой перпендикулярен \(AM\), если \(B,A,N\) коллинеарны?
  \item Какие два условия позволяют назвать \(AM\) серединным
  перпендикуляром?
  \item Что представляет собой пересечение двух серединных перпендикуляров?
  \item Почему перпендикуляр из \(M\) к \(BC\) проходит через середину \(BC\)?
  \item Где возникает прямоугольный равнобедренный треугольник?
  \item На какую дугу опирается \(\angle BCN\)?
  \item Из каких двух углов складывается \(\angle ABC\)?
\end{enumerate}

\begin{teacherbox}[Принцип выдачи подсказок]
Подсказку следует давать не в форме готового построения, а ступенчато:

\begin{enumerate}
  \item «Можно ли включить трапецию в более крупную фигуру?»
  \item «Что будет, если продолжить боковые стороны?»
  \item «Какую роль может играть отрезок, параллельный основанию и равный его
  половине?»
\end{enumerate}

Точка \(N\) называется только после того, как ученик исчерпал первые две
ступени.
\end{teacherbox}

\newpage
\section{Полное решение основной задачи}

\begin{answerbox}[Пункт а)]
Продлим стороны \(AB\) и \(CD\) до пересечения в точке \(N\). Получим
треугольник \(BNC\), причём \(AD\parallel BC\).

Треугольники \(NAD\) и \(NBC\) подобны. Так как
\[
\frac{AD}{BC}=\frac12,
\]
то
\[
\frac{NA}{NB}=\frac{ND}{NC}=\frac12.
\]
Следовательно, \(A\) и \(D\) — середины сторон \(BN\) и \(CN\), а \(AD\)
— средняя линия треугольника \(BNC\).

Поскольку \(B,A,N\) лежат на одной прямой и \(\angle BAM=90^\circ\),
имеем \(AM\perp BN\). Значит, \(AM\) — серединный перпендикуляр к \(BN\).
Аналогично \(DM\) — серединный перпендикуляр к \(CN\).

Итак, \(M\) — точка пересечения серединных перпендикуляров треугольника
\(BNC\), то есть центр его описанной окружности. Поэтому
\[
MB=MN=MC.
\]
В частности,
\[
\boxed{BM=CM}.
\]
\end{answerbox}

\circumcenterfigure

\begin{answerbox}[Пункт б)]
Обозначим \(AD=a\). Тогда \(BC=2a\). Пусть \(ME\perp BC\).

Из пункта а) \(MB=MC\), поэтому \(M\) лежит на серединном перпендикуляре
к \(BC\). Следовательно,
\[
BE=EC=\frac12BC=a.
\]
По условию расстояние от \(M\) до прямой \(BC\) равно \(AD\), значит
\[
ME=a.
\]
Треугольник \(BME\) прямоугольный равнобедренный, поэтому
\[
\angle MBC=45^\circ. \tag{1}
\]

Точки \(C,D,N\) лежат на одной прямой, поэтому
\[
\angle BCN=\angle BCD=64^\circ.
\]
Угол \(\angle BCN\) — вписанный, а \(\angle BMN\) — центральный; оба
опираются на дугу \(BN\). Следовательно,
\[
\angle BMN=2\cdot64^\circ=128^\circ.
\]
Треугольник \(BMN\) равнобедренный, так как \(MB=MN\). Поэтому
\[
\angle MBN=\frac{180^\circ-128^\circ}{2}=26^\circ. \tag{2}
\]

Луч \(BM\) лежит внутри угла \(NBC\), а лучи \(BA\) и \(BN\) совпадают.
Из (1) и (2):
\[
\angle ABC=\angle NBC
=\angle NBM+\angle MBC
=26^\circ+45^\circ
=\boxed{71^\circ}.
\]
\end{answerbox}

\Needspace{20\baselineskip}
\section{Разбор подготовительных заданий}

\begin{teacherbox}[A–B. Средняя линия]
Из подобия \(\triangle NAD\sim\triangle NBC\) и отношения \(AD:BC=1:2\)
получаем \(NA:NB=ND:NC=1:2\). Отсюда
\[
NA=AB,
\qquad
ND=DC.
\]
Важно не ограничиваться фразой «\(AD\) параллельно \(BC\), поэтому это
средняя линия»: одной параллельности недостаточно.
\end{teacherbox}

\begin{teacherbox}[C–D. Серединные перпендикуляры]
Ожидаемые ответы:
\[
AM\text{ — серединный перпендикуляр к }BN,
\qquad
MB=MN;
\]
\[
DM\text{ — серединный перпендикуляр к }CN,
\qquad
MC=MN.
\]
Следовательно, \(MB=MC=MN\), а \(M\) — центр описанной окружности.
\end{teacherbox}

\begin{teacherbox}[E–G. Углы]
В задании E:
\[
BE=EC=\frac12BC=a,
\qquad ME=a,
\qquad \angle MBC=45^\circ.
\]
В задании F:
\[
\angle BMN=2\angle BCN=128^\circ.
\]
В задании G:
\[
\angle MBN=\frac{180^\circ-128^\circ}{2}=26^\circ.
\]
\end{teacherbox}

\section{Типичные ошибки и коррекция}

\begin{tabularx}{\linewidth}{@{}>{\bfseries}p{43mm}X@{}}
\toprule
Ошибка & Корректирующий вопрос\\
\midrule
\(AD\) объявляют средней линией только из параллельности & «Откуда следует,
что концы \(AD\) являются серединами сторон?»\\
\midrule
Любой перпендикуляр называют серединным & «В какой точке он пересекает
отрезок? Доказано ли, что это середина?»\\
\midrule
\(M\) объявляют центром окружности по рисунку & «Какие равенства расстояний
уже доказаны?»\\
\midrule
Пишут \(N,M,E\) на одной прямой & «Доказано ли, что \(MN\perp BC\)?»\\
\midrule
Удваивают не тот угол & «На какую дугу опирается каждый из выбранных углов?»\\
\midrule
Получают \(26^\circ\), но не могут собрать ответ & «Какой луч делит угол
\(ABC\) на две найденные части?»\\
\bottomrule
\end{tabularx}

\begin{warningbox}[О единственности перпендикуляра]
Единственность перпендикуляра можно применить только так: если отдельно
доказано, что \(ME\perp BC\) и \(MN\perp BC\), то прямые \(ME\) и \(MN\),
проходящие через одну точку \(M\), совпадают. В задаче второе утверждение
неверно, поэтому точка \(N\) не лежит на \(ME\).
\end{warningbox}

\newpage
\section{Ответы к домашней работе}

\begin{answerbox}[Задание 1]
Так как \(K\) — середина \(PQ\) и \(OK\perp PQ\), прямая \(OK\) является
серединным перпендикуляром к \(PQ\). Следовательно, \(OP=OQ\).
\end{answerbox}

\begin{answerbox}[Задание 2]
Из принадлежности \(O\) серединному перпендикуляру к \(XY\) следует
\(OX=OY\), а из принадлежности серединному перпендикуляру к \(XZ\) —
\(OX=OZ\). Поэтому \(OY=OZ\). Точка \(O\) — центр описанной окружности
треугольника \(XYZ\).
\end{answerbox}

\begin{answerbox}[Задание 3]
\[
\angle AOB=2\angle ACB=114^\circ.
\]
Поскольку \(OA=OB\), углы при основании треугольника \(AOB\) равны
\[
\frac{180^\circ-114^\circ}{2}=33^\circ.
\]
\end{answerbox}

\begin{answerbox}[Задание 4]
Перпендикуляр из центра окружности к хорде делит её пополам, поэтому
\(BE=6\) см. По условию \(OE=6\) см. Треугольник \(OBE\) прямоугольный
равнобедренный, следовательно,
\[
\angle OBC=45^\circ.
\]
\end{answerbox}

\Needspace{18\baselineskip}
\begin{answerbox}[Задание 5]
Пусть \(E\) — основание перпендикуляра из \(M\) к \(BC\). Тогда
\(BE=\frac12BC=ME\), поэтому \(\angle MBC=45^\circ\).

Далее
\[
\angle BMN=2\angle BCN=124^\circ,
\]
а из \(MB=MN\)
\[
\angle MBN=\frac{180^\circ-124^\circ}{2}=28^\circ.
\]
Следовательно,
\[
\angle NBC=28^\circ+45^\circ=\boxed{73^\circ}.
\]
\end{answerbox}

\begin{answerbox}[Задание 6]
Решение полностью повторяет маршрут основной задачи. Получаем
\[
\angle MBC=45^\circ,
\qquad
\angle MBN=90^\circ-68^\circ=22^\circ.
\]
Поэтому
\[
\angle ABC=45^\circ+22^\circ=\boxed{67^\circ}.
\]
\end{answerbox}

\begin{answerbox}[Задание 7]
Единственность перпендикуляра не доказывает принадлежность точки \(N\)
прямой \(ME\). Для такого вывода нужно было бы отдельно доказать
\[
MN\perp BC.
\]
Только тогда прямые \(MN\) и \(ME\) совпали бы. В данной конфигурации
\(MN\perp BC\) вообще не выполняется.
\end{answerbox}

\Needspace{12\baselineskip}
\begin{answerbox}[Задание 8]
Угол \(\angle MBC=45^\circ\). Центральный угол
\[
\angle BMN=2\alpha,
\]
поэтому в равнобедренном треугольнике \(BMN\)
\[
\angle MBN=\frac{180^\circ-2\alpha}{2}=90^\circ-\alpha.
\]
Следовательно,
\[
\boxed{\angle ABC=135^\circ-\alpha}.
\]
\end{answerbox}

\section{Критерии оценки самостоятельного решения}

\begin{tabularx}{\linewidth}{@{}p{16mm}X@{}}
\toprule
Баллы & Критерий\\
\midrule
1 & Выполнено достроение до треугольника \(BNC\).\\
1 & Обосновано, что \(A,D\) — середины сторон.\\
2 & Корректно распознаны два серединных перпендикуляра.\\
1 & Сделан вывод \(M\) — центр описанной окружности и доказано \(BM=CM\).\\
2 & Обосновано получение \(\angle MBC=45^\circ\).\\
2 & Корректно получены \(128^\circ\) и \(26^\circ\).\\
1 & Правильно сложены углы и записан ответ \(71^\circ\).\\
\bottomrule
\end{tabularx}

\begin{teacherbox}[Уровни освоения]
\begin{itemize}
  \item \textbf{Минимальный:} ученик с подсказкой выполняет достроение и
  решает пункт а).
  \item \textbf{Базовый:} самостоятельно решает пункт а), а пункт б) — по
  направляющим вопросам.
  \item \textbf{Повышенный:} самостоятельно записывает полное решение.
  \item \textbf{Высокий:} получает и обосновывает формулу
  \(\angle ABC=135^\circ-\alpha\).
\end{itemize}
\end{teacherbox}
